\documentclass[11pt,a4paper]{article}
\usepackage[margin=0.85in,top=1in,bottom=1in]{geometry}
\usepackage{amsmath,amssymb,amsfonts}
\usepackage{graphicx}
\usepackage{float}
\usepackage{enumitem}
\usepackage{microtype}
\usepackage{caption}
\usepackage[hidelinks]{hyperref}
\usepackage[most,breakable]{tcolorbox}
\tcbuselibrary{breakable,skins}
\usepackage{fontspec}
\setmainfont{DejaVu Serif}
\setsansfont{DejaVu Sans}
\setmonofont{DejaVu Sans Mono}
\captionsetup{font=small,labelfont=bf,skip=4pt}
\setlist{leftmargin=1.5em,topsep=2pt}

%% Breakable card — prevents content cutoff across pages
\newtcolorbox{solcard}[3]{
  breakable, enhanced,
  colback=white, colframe=gray!50, arc=2pt,
  title={\textbf{#1}\hfill\textsf{\small #3\,pts}},
  fonttitle=\normalsize\sffamily\bfseries,
  before upper={\small\textit{#2}\par\smallskip\hrule height 0.4pt\smallskip},
  top=4pt, bottom=6pt, left=7pt, right=7pt,
  parbox=false,
}

\title{\textbf{In the footsteps of van der Waals --- Solution}}
\author{\normalsize W22 (2022) — English translation}
\date{}
\begin{document}\maketitle\vspace{-0.5em}
\begin{solcard}{A1}{Write down an expression for finding the optimal value $a_0$ through the sums (or averages) of values ​​$T, V$ and their combinations.}{0.40}

We will solve the problem formally, express the temperature from the model function and write down the sum $S$ that needs to be minimized. $$T=\frac{\left(p+a_{0}\right) V}{R}$$ $$S=\sum_{n}\left(T_{i}-\frac{p V_i}{R}-\frac{a_{0} V_{i}}{R}\right)^{2}$$ Differentiate: $$\frac{\partial S}{\partial a_{0}}=\sum_{n} 2\left(T_{i}-\frac{\left(p+a_{0}\right) V_{i}}{R}\right)\left(-\frac{V_{i}}{R}\right)=0$$ Transform (сократив on "-2"): $$\frac{1}{R} \sum_{n} T_{i} V_{i}-\frac{p+a_{0}}{R^{2}} \sum_{n} V_{i}^{2}=0$$ and further $$p+a_{0}=\frac{R \sum_{n} T_{i} V_{i}}{\sum_{n} V_{i}^{2}}$$ We finally get $$a_{0}=\frac{R \sum_{n} T_{i} V_{i}}{\sum_n V_{i}^{2}}-p=R \frac{\overline{T V}}{\overline{V^{2}}}-p$$ $\textit{Note:}$ The answer could be written almost immediately if you realize that the model formula is linear dependence with one parameter: slope.

\end{solcard}

\begin{solcard}{A2}{Calculate the sums (or averages) you need for the entire data set.}{0.30}

$\sum_{n} T_{i} V_{i}=6.5903~\mathrm{K}\cdot \mathrm{m}^3, \overline{T V}=0.31382~\mathrm{K}\cdot \mathrm{m}^3$ $\sum_n V_{i}^{2}=1.8935\cdot 10^{-5}~\mathrm{m}^6, \overline{V^{2}}=9.0167\cdot10^{-7}~\mathrm{m}^6$ Here and below, we save more significant figures to reduce the rounding error of the final answer.

\end{solcard}

\begin{solcard}{A3}{Calculate the value of $a_0$. Specify it to three significant figures.}{0.40}

$a_0=892~\mathrm{kPa}$

\end{solcard}

\begin{solcard}{A4}{For each volume value, calculate the temperature value that is obtained as a result of fitting. Plot the initial experimental data and fitting points on one graph. Mark data from different series with different icons. Sign the series.}{0.40}
\begin{figure}[H]
  \centering
  \includegraphics[width=0.72\textwidth]{PE_S_EN_assets/fig1.jpg}
\end{figure}
\end{solcard}

\begin{solcard}{B1}{Write down an expression for finding the optimal value $a_1$ through the sums (or averages) of values ​​$T, V$ and their combinations.}{0.40}

Again, we act formally: $$T=\frac{p V}{R}+\frac{a_{1}}{R}$$ $$S=\sum_{n}\left(T_{i}-\frac{p V_{i}}{R}-\frac{a_{1}}{R}\right)^{2}$$ $$\frac{\partial S}{\partial a_{1}}=\sum_{n} 2\left(T_{i}-\frac{p V_{i}}{R}-\frac{a_{1}}{R}\right)\left(-\frac{1}{R}\right)=0$$ $$\sum_{n} T_{i}-\frac{p}{R} \sum_{n} V_{i}-\frac{a_{1}}{R} \sum_{n} 1=0$$ Hence we get окончательный answer: $$a_{1}=\frac{R \sum_{n} T_{i}-p \sum_{n} V_{i}}{n}=R \overline{T}-p \overline{V}$$

\end{solcard}

\begin{solcard}{B2}{Calculate the sums (or averages) you need for the entire data set.}{0.30}

$\sum_{n} T_{i}=7009.2~\mathrm{K}, \overline{T}=333.771~\mathrm{K}$ $\sum_{n} V_{i}=0.0189~\mathrm{m}^3, \overline{V}=0.0009~\mathrm{m}^3$

\end{solcard}

\begin{solcard}{B3}{Calculate the value of $a_1$. Specify it to three significant figures.}{0.40}

$a_1=974~\mathrm{Pa}\cdot \mathrm{m}^3$

\end{solcard}

\begin{solcard}{B4}{For each volume value, calculate the temperature value that is obtained as a result of fitting. Plot the initial experimental data and fitting points on one graph. Mark data from different series with different icons. Sign the series.}{0.40}
\begin{figure}[H]
  \centering
  \includegraphics[width=0.72\textwidth]{PE_S_EN_assets/fig2.jpg}
\end{figure}
\end{solcard}

\begin{solcard}{C1}{Write down an expression for finding the optimal value $a_2$ through the sums (or averages) of values ​​$T, V$ and their combinations.}{0.40}

Again, we act formally: $$T=\frac{p V}{R}+\frac{a_{2}}{V R}$$ $$S=\sum_{n}\left(T_{i}-\frac{p V_{i}}{R}-\frac{a_{2}}{V_{i} R}\right)^{2}$$ $$\frac{\partial S}{\partial a_{2}}=\sum_{n} 2\left(T_{i}-\frac{p V_{i}}{R}-\frac{a_{2}}{V_{i} R}\right)\left(-\frac{1}{V_{i} R}\right)=0$$ $$-\frac{1}{R} \sum_{n} \frac{T_{i}}{V_{i}}+\frac{p}{R^{2}} \sum_{n} 1+\frac{a_{2}}{R^{2}} \sum_{n} \frac{1}{V_{i}^{2}}=0$$ Hence we finally get: $$a_{2}=\frac{R \sum_{n} \frac{T_{i}}{V_{i}}-p n}{\sum_{n} \frac{1}{V_{i}^{2}}}=\frac{R\overline{\left(\frac{T}{V}\right)}-p}{\overline{\left(\frac{1}{V^{2}}\right)}}$$ $\textit{Note:}$ Интуитивно one can be бы пойти путем линеаризации $p V^{2}+a_{2}=R T V$, and by аналогии with part B1 write, that $a_{2}=R \overline{T V}-p \overline{V^{2}}$. This erroneous path and incorrect solution. For calculation such method estimate for $a_2$ will be “displaced”. In C3, we will calculate the values ​​using both methods and see that using linearization results in a higher sum of squared deviations.

\end{solcard}

\begin{solcard}{C2}{Calculate the sums (or averages) you need for the entire data set.}{0.30}

$\sum_{n} \frac{T_{i}}{V_{i}}=8489686~\mathrm{K}/м^3, \overline{\left(\frac{T}{V}\right)}=404271~\mathrm{K}/м^3$ $\sum_{n} \frac{1}{V_{i}^{2}}= 39221160~\mathrm{m}^{-6}, \overline{\left(\frac{1}{V^{2}}\right)}=1867674~\mathrm{m}^{-6}$

\end{solcard}

\begin{solcard}{C3}{Calculate the value of $a_2$. Specify it to three significant figures.}{0.40}

$a_2=0.728~\mathrm{Pa}\cdot\mathrm{m}^6$ To calculate $a_2^*$ for the linearized dependence, we have already calculated all the quantities (point A2). Therefore $a_2^*=0.805$. If you calculate $S(a_2^*)=8813~\mathrm{K}^2$, it will be greater than $S(a_2)=5478~\mathrm{K}^2$.

\end{solcard}

\begin{solcard}{C4}{For each volume value, calculate the temperature value that is obtained as a result of fitting. Plot the initial experimental data and fitting points on one graph. Mark data from different series with different icons. Sign the series.}{0.40}
\begin{figure}[H]
  \centering
  \includegraphics[width=0.72\textwidth]{PE_S_EN_assets/fig3.jpg}
\end{figure}
\end{solcard}

\begin{solcard}{D1}{Write down expressions for finding the optimal values ​​of $c$ and $d$ through the sums (or averages) of values ​​$T, V$ and their combinations.}{1.50}

Let's act again formally: $$T=\frac{p V}{R}+\frac{C}{R V}+\frac{d}{R V^{2}}$$ $$S=\sum_{n}\left(T_{i}-\frac{p V_{i}}{R}-\frac{c}{R V_{i}}-\frac{d}{R V_{i}^{2}}\right)^{2}$$ Differentiate, we get the system of equations: $$\frac{\partial S}{\partial c}=\sum_{n} 2\left(T_{i}-\frac{p V_{i}}{R}-\frac{c}{R V_{i}}-\frac{d}{R V_{i}^{2}}\right)\left(-\frac{1}{R V_{i}}\right)=0$$ $$\frac{\partial S}{\partial d}=\sum_{n} 2\left(T_{i}-\frac{p V_{i}}{R}-\frac{C}{R V_{i}}-\frac{d}{R V_{i}^{2}}\right)\left(-\frac{1}{R V_{i}^{2}}\right)=0$$ Transform system: $$R \sum_{n} \frac{T_{i}}{V_{i}}-p \sum_{n} 1-c \sum_{n} \frac{1}{V_{i}^{2}}-d \sum_{n} \frac{1}{V_{i}^{3}}=0$$ $$R \sum_{n} \frac{T_{i}}{V_{i}^{2}}-p \sum_{n} \frac{1}{V_{i}}-c \sum_{n} \frac{1}{V_{i}^{3}}-d \sum_{n} \frac{1}{V_{i}^{4}}=0$$ Transition to mean value and lead to system to convenient form: $$c \overline{V^{-2}}+d \overline{V^{-3}}=R \overline{T / V}-p$$ $⟦ 14⟧$ From which, solve jointly equation, obtain answer: $$c=\frac{R\left(\overline{T / V^{2}}\,\overline{V^{-3}}-\overline{T / V}\,\overline{V^{-4}}\right)-p\left(\overline{V^{-1}}\,\overline{V^{-3}}-\overline{V^{-4}}\right)}{\left(\overline{V^{-3}}\right)^{2}-\overline{V^{-2}}\,\overline{V^{-4}}}$$ $$d=\frac{R\left(\overline{T / V^{2}}\,\overline{V^{-2}}-\overline{T / V}\,\overline{V^{-3}}\right)-p\left(\overline{V^{-1}}\,\overline{V^{-2}}-\overline{V^{-3}}\right)}{\overline{V^{-2}}\,\overline{V^{-4}}-\left(\overline{V^{-3}}\right)^{2}}$$ $\textit{Note:}$ Similarly part C1, one can be бы lead to модельную function to linear dependence: $c V+d=R T V^{2}-p V^{3}$, and determine coefficient $c$ and $d$ by formula, lead to in introduction to problem. This erroneous path and incorrect solution. For calculation such method estimate for $c$ and $d$ will be “displaced”.

\end{solcard}

\begin{solcard}{D2}{Calculate the sums (or averages) you need for the entire data set.}{2.00}

We calculated two values ​​$\overline{T / V}$ and $\overline{V^{-2}}$ in point C2. $\overline{V^{-1}}=1270.76~\mathrm{m}^{-3}$ $\overline{V^{-3}}=3.13694\cdot 10^9~\mathrm{m}^{-9}$ $\overline{V^{-4}}=5.85765\cdot 10^{12}~\mathrm{m}^{-12}$ $\overline{T / V^2}=0.568387\cdot 10^9~\mathrm{K}/м^6$

\end{solcard}

\begin{solcard}{D3}{Calculate the values ​​of $c$ and $d$. Please provide them to three significant figures.}{1.00}

$c=1.02~\mathrm{Pa}\cdot\mathrm{m}^6$ $d=-1.73\cdot10^{-4}~\mathrm{Pa}\cdot\mathrm{m}^9$ Offset values: $c^*=1.08~\mathrm{Pa}\cdot\mathrm{m}^6, d^*=-2.17\cdot10^{-4}~\mathrm{Pa}\cdot\mathrm{m}^9$

\end{solcard}

\begin{solcard}{D4}{For each volume value, calculate the temperature value that is obtained as a result of fitting. Plot the initial experimental data and fitting points on one graph. Mark data from different series with different icons. Sign the series.}{1.00}
\begin{figure}[H]
  \centering
  \includegraphics[width=0.72\textwidth]{PE_S_EN_assets/fig4.jpg}
\end{figure}
\end{solcard}

\begin{solcard}{E1}{Calculate the values ​​of $a$ and $b$. Please provide them to three significant figures. Consider the scale of work carried out by Van der Waals.}{??}

In this part of the problem, it is necessary to similarly formally minimize the sum over the parameters $a$ and $b$. The final system of equations for the parameters will look like this: 
\[
\begin{cases} \alpha_1 a + \alpha_2 b + \alpha_3 ab + \alpha_4 a^2 + \alpha_5 a^2b = c_1\\ \beta_1 a + \beta_2 b + \beta_3 ab + \beta_4 b^2 + \beta_5 ab^2 = c_2 \end{cases}
\] Coefficients $\alpha_i$ and $\beta_i$ are related entirely to the average values calculated in point D2. Those. There would be no additional need to calculate anything. Each of the equations is linear in one of the variables and quadratic in the second. Therefore, the system can be reduced to a fifth-degree equation for one of the variables. An equation of degree five definitely has at least one valid solution.

\end{solcard}

\end{document}